\spacing{1.3}


\section{Sistemas materiales}
Podemos clasificar los sistemas materiales en (ver esquema):
\vspace{12pt}

\esquema{Sistemas materiales}
{\esquema{Sustancias puras}
{\esquema{Elementos}
{\( \bullet \) Hidrógeno
\vspace{6pt}\\
    \( \bullet \) Litio}
\vspace{20pt}\\
\esquema{Compuestos}
{\( \bullet \)\ch{NaCl}
\vspace{6pt}\\
    \( \bullet \)\ch{H2O}
\vspace{6pt}\\
    \( \bullet \)\ch{H2SO4}}
}
\vspace{20pt} \\
\esquema{Mezclas}
{\( \bullet \) Homogéneas
\vspace{6pt}\\
    \( \bullet \) Heterogéneas
}
}

\begin{enumerate}[label=\numberballmDarkTeal]
\item \negsubro{Sustancias puras}: un tipo de materia que tiene una composición constante y propiedades definidas, y que no puede separarse en otras sustancias mediante métodos físicos.
Las sustancias puras pueden dividirse en:
\begin{itemize}
    \item \textbf{Elementos}: formados por un único tipo de átomo.
    \item \textbf{Compuestos}: formados por la unión de dos o más elementos químicos.
\end{itemize}

\item \negsubro{Mezclas}: la combinación de dos o más sustancias puras que pueden separarse mediante métodos físicos.
Las mezclas pueden dividirse en:
\begin{itemize}
    \item \textbf{Mezcla homogénea}: una mezcla en la que los componentes están distribuidos de manera uniforme y no pueden distinguirse a simple vista.
    \item \textbf{Mezcla heterogénea}: una mezcla en la que los componentes pueden distinguirse.
\end{itemize}

\end{enumerate}


\section{Disoluciones}

\begin{definition}[Disolución]
    Es un sistema homogéneo formado por dos o más sustancias puras en distintas proporciones.
\end{definition}

Los componentes de una disolución son:
\begin{itemize}
    \item \textbf{Disolvente}. La sustancia en mayor cantidad.
    \item \textbf{Soluto}. La sustancia en menor cantidad. Puede haber más de un soluto.
    \item \textbf{Disolución}. La mezcla del soluto(s) y el disolvente.
\end{itemize}

Por ejemplo, en una disolución de agua y azúcar, el disolvente es el agua y el soluto es el azúcar.
En una botella de agua mineral, la etiqueta indica que contiene sodio, calcio, magnesio, etc.
En este caso, el disolvente es el agua, y los solutos son el sodio, el calcio y el magnesio.

\subsection{Concentración de una disolución}

Cuando tenemos una disolución, es fundamental saber cuánto soluto hay en el disolvente.

\begin{definition}[Concentración]
    Es la cantidad de soluto disuelto en una cierta cantidad de disolvente.
\end{definition}

Existen diferentes formas de expresar la concentración. Las principales son:

\begin{enumerate}[label=\numberballmDarkTeal]
\item \negsubro{Porcentaje en masa}.
Nos indica la masa de soluto en cada \qty{100}{\g} de disolución.

\begin{equation}
    \text{\% en masa de soluto} = \dfrac{\text{masa de soluto}}{\text{masa de disolución}}\cdot 100
    \label{eq:porcentaje_masa}
\end{equation}

Importante:
\begin{itemize}
    \item Si hay varios solutos, calculamos el porcentaje de cada uno.
    \item También podemos calcular el porcentaje del disolvente.
    \item Todos los porcentajes juntos deben sumar \num{100}.
\end{itemize}
\end{enumerate}

\begin{example}
    Tenemos una disolución acuosa con sodio y calcio.
    Hay \qty{18}{\g} de sodio, \qty{6}{\g} de calcio, y \qty{100}{\g} de agua.
    Calcula el porcentaje en masa del sodio y del calcio.

    En ambos casos debemos aplicar la ecuación \eqref{eq:porcentaje_masa}.
    En ambos casos tenemos que:

    \begin{equation*}
        \text{\% masa de sodio} = \dfrac{\qty{18}{\g}\text{ de sodio}}{\qty{18}{\g}\text{ de sodio} + \qty{6}{\g}\text{ de calcio} + \qty{100}{\g}\text{ de agua}}\cdot 100 = 14.52 \%
    \end{equation*}

    \begin{equation*}
        \text{\% masa de calcio} = \dfrac{\qty{6}{\g}\text{ de calcio}}{\qty{18}{\g}\text{ de sodio} + \qty{6}{\g}\text{ de calcio} + \qty{100}{\g}\text{ de agua}}\cdot 100 = 4.84 \%
    \end{equation*}

    Si también queremos calcular el porcentaje del disolvente, sería:

    \begin{equation*}
        \text{\% masa de agua} = \dfrac{\qty{100}{\g}\text{ de agua}}{\qty{18}{\g}\text{ de sodio} + \qty{6}{\g}\text{ de calcio} + \qty{100}{\g}\text{ de agua}}\cdot 100 = 80.65 \%
    \end{equation*}
\end{example}

\begin{exercise}
    El acero es una aleación compuesta de hierro, carbono y otros elementos. Una masa de \qty{120}{\g} de acero contiene un 1,8\% de carbono y un 1,2\% de manganeso. Calcula la masa de cada elemento.
    \begin{solu}
        \textbf{Fórmula:} $m_{\text{componente}} = \frac{\% \text{ en masa}}{100} \cdot m_{\text{total}}$ \\
        \textbf{Datos:} $m_{\text{total}} = 120 \text{ g}$; $\%C = 1.8\%$; $\%Mn = 1.2\%$ \\
        \textbf{Solución:}
        \begin{enumerate}
            \item \textbf{Masa de carbono:} $m_{C} = \frac{1.8}{100} \cdot 120 \text{ g} = 2.16 \text{ g}$
            \item \textbf{Masa de manganeso:} $m_{Mn} = \frac{1.2}{100} \cdot 120 \text{ g} = 1.44 \text{ g}$
            \item \textbf{Masa de hierro:} Se calcula restando los demás componentes del total: \\
            $m_{Fe} = 120 \text{ g} - (2.16 \text{ g} + 1.44 \text{ g}) = 116.4 \text{ g}$
        \end{enumerate}
    \end{solu}
\end{exercise}

\begin{exercise}
    Disolvemos 12 g de cloruro de sodio y 13 g de cloruro de potasio en 250 g de agua destilada. Calcula el porcentaje en masa de cada soluto en la disolución final.
    \begin{solu}
        \textbf{Fórmulas:} $m_{\text{total}} = m_{s1} + m_{s2} + m_{\text{disolvente}}$ y $\%m = \frac{m_{\text{soluto}}}{m_{\text{total}}} \cdot 100$ \\
        \textbf{Solución:}
        \begin{enumerate}
            \item \textbf{Masa total:} $12 \text{ g} + 13 \text{ g} + 250 \text{ g} = 275 \text{ g}$.
            \item \textbf{\% NaCl:} $\frac{12 \text{ g}}{275 \text{ g}} \cdot 100 \approx 4.36\%$
            \item \textbf{\% KCl:} $\frac{13 \text{ g}}{275 \text{ g}} \cdot 100 \approx 4.73\%$
        \end{enumerate}
    \end{solu}
\end{exercise}

\begin{exercise}
    Queremos preparar 500 g de una disolución acuosa que contenga un 0,5\% de cloruro de sodio, un 0,3\% de cloruro de potasio y un 0,2\% de nitrato de potasio en masa. ¿Cuánto de cada soluto se necesita? ¿Cuánta agua necesitamos?
    \begin{solu}
        \textbf{Fórmula:} $m_s = \frac{\%m}{100} \cdot m_{\text{total}}$ \\

        \textbf{Datos:} $m_{\text{total}} = 500 \text{ g}$. \\

        \textbf{Solución:}
        \begin{itemize}
            \item $m_{NaCl} = 500 \cdot 0.005 = 2.5 \text{ g}$
            \item $m_{KCl} = 500 \cdot 0.003 = 1.5 \text{ g}$
            \item $m_{KNO_3} = 500 \cdot 0.002 = 1.0 \text{ g}$
            \item \textbf{Masa de agua:} $500 \text{ g} - (2.5 + 1.5 + 1.0) \text{ g} = 495 \text{ g}$
        \end{itemize}
    \end{solu}
\end{exercise}

\begin{exercise}
    ¿Qué significa que una disolución de yodo en alcohol sea del 1\% en masa?
    \begin{solu}
        \textbf{Definición de porcentaje en masa:} Representa las partes de soluto en 100 partes de disolución. \\
        Significa que en cada \textbf{100 g de disolución}, hay exactamente \textbf{1 g de yodo} (soluto) y los \textbf{99 g restantes son alcohol} (disolvente).
    \end{solu}
\end{exercise}

\begin{exercise}
    Utilizamos una mezcla de plomo y estaño en soldadura. Si el porcentaje en masa de plomo es del 66\% y el de estaño del 34\%, ¿qué masas de cada metal fundirías para preparar 250 g de esta aleación?
    \begin{solu}
        \textbf{Fórmula:} $m_{\text{componente}} = \frac{\%m}{100} \cdot m_{\text{total}}$ \\

        \textbf{Datos:} $m_{\text{total}} = 250 \text{ g}$. \\

        \textbf{Solución:}
        \begin{itemize}
            \item Masa de plomo (Pb): $250 \cdot 0.66 = 165 \text{ g}$
            \item Masa de estaño (Sn): $250 \cdot 0.34 = 85 \text{ g}$
        \end{itemize}
        (Comprobación: $165 + 85 = 250 \text{ g}$).
    \end{solu}
\end{exercise}


\begin{exercise}
    El porcentaje en masa de cloruro de potasio en una cierta disolución es del 10\%. ¿Cuánto de este soluto está disuelto en 250 g de disolución?
    \begin{solu}
        \textbf{Fórmula:} $m_s = \frac{\%m}{100} \cdot m_{\text{disolución}}$ \\
        \textbf{Solución:} Aplicamos la definición directamente a la masa total: \\
        $m_{KCl} = 250 \text{ g} \cdot \frac{10}{100} = 25 \text{ g}$ de KCl.
    \end{solu}
\end{exercise}

\begin{exercise}
    Un agua embotellada contiene un 99,8\% en masa de agua y un 0,2\% de dióxido de carbono. Identifica el soluto y el disolvente en esta disolución. ¿Qué masa de dióxido de carbono contendría 1 kg de agua?
    \begin{solu}
        \textbf{Concepto:} Identificar el soluto (menor cantidad) y el disolvente (mayor cantidad). \\
        \textbf{Fórmula de proporción:} $\frac{m_{\text{soluto}}}{m_{\text{disolvente}}} = \frac{\% \text{ soluto}}{\% \text{ disolvente}}$ \\
        \textbf{Solución:}
        \begin{itemize}
            \item \textbf{Identificación:} Soluto = dióxido de carbono ($CO_2$); Disolvente = agua ($H_2O$).
            \item \textbf{Cálculo de la masa:} Tenemos 1 kg (1000 g) de agua. \\
            $m_{CO_2} = \frac{0.2 \text{ g } CO_2}{99.8 \text{ g } H_2O} \cdot 1000 \text{ g } H_2O \approx 2.004 \text{ g}$ de $CO_2$.
        \end{itemize}
    \end{solu}
\end{exercise}

\begin{enumerate}[label=\numberballmDarkTeal,resume]
\item \negsubro{Porcentaje en volumen}.
Nos indica el volumen de soluto en cada \qty{100}{\litre} de disolución.

\begin{equation}
    \text{\% en volumen de soluto} = \dfrac{\text{volumen de soluto}}{\text{volumen de disolución}}\cdot 100
\end{equation}
\end{enumerate}

\begin{example}
    Preparamos una disolución con \qty{100}{\ml} de alcohol, \qty{150}{\ml} de acetona, y \qty{250}{\ml} de agua.
    Calcula el porcentaje en volumen de cada sustancia.

    \begin{equation*}
        \text{\% vol alcohol} = \dfrac{\qty{100}{\ml}\text{ de alcohol}}{\qty{100}{\ml}\text{ de alcohol} + \qty{150}{\ml}\text{ de acetona} + \qty{250}{\ml}\text{ de agua}}\cdot 100 = 20 \%
    \end{equation*}

    \begin{equation*}
        \text{\% vol acetona} = \dfrac{\qty{150}{\ml}\text{ de acetona}}{\qty{100}{\ml}\text{ de alcohol} + \qty{150}{\ml}\text{ de acetona} + \qty{250}{\ml}\text{ de agua}}\cdot 100 = 30 \%
    \end{equation*}

    \begin{equation*}
        \text{\% vol agua} = \dfrac{\qty{250}{\ml}\text{ de agua}}{\qty{100}{\ml}\text{ de alcohol} + \qty{150}{\ml}\text{ de acetona} + \qty{250}{\ml}\text{ de agua}}\cdot 100 = 50 \%
    \end{equation*}
\end{example}

\begin{exercise}
    Se prepara una disolución de éter y cloroformo añadiendo \qty{10}{\mL} de éter a 90 mL de cloroformo. ¿Cuál es el porcentaje en volumen de cada componente?
    \begin{solu}
        \textbf{Fórmulas:} \\
        1. $V_{\text{disolución}} = V_{\text{soluto}} + V_{\text{disolvente}}$ \\
        2. $\%V = \frac{V_{\text{componente}}}{V_{\text{disolución}}} \cdot 100$ \\

        \textbf{Datos:} $V_{\text{éter}} = 10 \text{ mL}$; $V_{\text{cloroformo}} = 90 \text{ mL}$ \\

        \textbf{Solución:}
        \begin{enumerate}
            \item \textbf{Volumen de la disolución:} $V_{\text{total}} = 10 \text{ mL} + 90 \text{ mL} = 100 \text{ mL}$.
            \item \textbf{Porcentaje de éter:} $\%V_{\text{éter}} = \frac{10 \text{ mL}}{100 \text{ mL}} \cdot 100 = 10\%$
            \item \textbf{Porcentaje de cloroformo:} $\%V_{\text{cloro}} = \frac{90 \text{ mL}}{100 \text{ mL}} \cdot 100 = 90\%$
        \end{enumerate}
    \end{solu}
\end{exercise}

\begin{exercise}
    La composición en volumen de los ingredientes activos en un envase de 200 mL de loción corporal es: 2,14\% de aceite de almendras dulces; 2,14\% de aceite de maíz; 4,29\% de aceite de cacahuete y 2,88\% de aceite de copra hidrogenado. Calcula el volumen de cada componente.
    \begin{solu}
        \textbf{Fórmula:} $V_{\text{comp}} = \frac{\%V}{100} \cdot V_{\text{total}}$ \\
        \textbf{Datos:} $V_{\text{total}} = 200 \text{ mL}$. \\
        \textbf{Solución:} Multiplicamos el volumen total por cada porcentaje decimal:
        \begin{itemize}
            \item Aceite de almendras: $200 \cdot 0.0214 = 4.28 \text{ mL}$
            \item Aceite de maíz: $200 \cdot 0.0214 = 4.28 \text{ mL}$
            \item Aceite de cacahuete: $200 \cdot 0.0429 = 8.58 \text{ mL}$
            \item Aceite de copra: $200 \cdot 0.0288 = 5.76 \text{ mL}$
        \end{itemize}
    \end{solu}
\end{exercise}


\begin{exercise}
    Hemos disuelto \qty{25}{\cmcubicos} de alcohol en agua destilada para obtener un volumen total de disolución de \qty{250}{\cmcubicos}.
    \begin{enumerate}[label=\letraminuscula]
        \item ¿Cuál es el volumen del soluto y del disolvente?
        \item Calcula el porcentaje en volumen de cada componente de esta disolución.
    \end{enumerate}
    \begin{solu}
        \textbf{Fórmulas:} $V_{\text{disolvente}} = V_{\text{total}} - V_{\text{soluto}}$ y $\%V = \frac{V_{\text{comp}}}{V_{\text{total}}} \cdot 100$ \\
        \textbf{Solución:}
        \begin{enumerate}[label=\alph*)]
            \item \textbf{Volúmenes:} Soluto = $25 \text{ cm}^3$ (alcohol). \\
            Disolvente = $250 \text{ cm}^3 - 25 \text{ cm}^3 = 225 \text{ cm}^3$ (agua).
            \item \textbf{Porcentajes:} \\
            $\%V_{\text{alcohol}} = \frac{25}{250} \cdot 100 = 10\%$ \\
            $\%V_{\text{agua}} = \frac{225}{250} \cdot 100 = 90\%$
        \end{enumerate}
    \end{solu}
\end{exercise}


\begin{enumerate}[label=\numberballmDarkTeal,resume]
\item \negsubro{Concentración en masa}.
Nos indica la masa de soluto por unidad de volumen de disolución.
Se suele expresar en \unit{\g\per\litre}.

\begin{equation*}
    \text{Concentración en masa} = \dfrac{m_{\text{soluto (g)}}}{V_{\text{disolución (L)}}}
\end{equation*}

No hay que confundirla con la densidad:
\begin{itemize}
    \item La densidad es la masa de una sustancia por unidad de volumen.
    \item La concentración en masa es la masa de soluto por unidad de volumen de la \textbf{disolución}.
\end{itemize}
\end{enumerate}

\begin{example}
    Para preparar una disolución de nitrato de potasio con \qty{20}{\g} de nitrato en agua para obtener un volumen final de disolución de \qty{200}{\ml}. Calcula:
    \begin{enumerate}[label=\letraminuscula]
\item La concentración en masa de la disolución expresada en \unit{\g\per\litre}.

Como el volumen de la disolución se da en \unit{\ml}, lo convertimos a \unit{\litre}. El volumen es entonces \qty{0.2}{\litre}. Si aplicamos:
\begin{equation*}
    \text{Concentración en masa} = \dfrac{m_{\text{soluto (g)}}}{V_{\text{disolución (L)}}} = \dfrac{20}{0.2} = \qty{100}{\g\per\litre}
\end{equation*}

\item Calcula la masa de la disolución si la densidad de esta disolución es \( \rho = \qty{1.09}{\g\per\cubic\centi\metre} \).

Para calcular su masa, despejamos la masa en la ecuación de la densidad:
\begin{equation*}
    \rho = \dfrac{m}{V} \rightarrow m = \rho \cdot V
\end{equation*}
Antes de sustituir en la ecuación, debemos comprobar que tanto la densidad como el volumen están en las mismas unidades, pero sabemos que \( \qty{1}{\ml} = \qty{1}{\cubic\centi\metre} \).
\begin{equation*}
    m = \rho \cdot V = 1.09 \cdot 200 = \qty{218}{\g}
\end{equation*}
    \end{enumerate}
\end{example}

\begin{exercise}
    200 ml de una disolución acuosa contienen 12 g de azúcar. La densidad de esta disolución es \qty{1.022}{\g\per\cmcubicos}.
    \begin{enumerate}[label=\letraminuscula]
        \item Calcula el porcentaje en masa del azúcar.
        \item Calcula la concentración en masa en g/L.
    \end{enumerate}
    \begin{solu}
        \textbf{Fórmulas:} $m_{\text{disolución}} = d \cdot V$; $\%m = \frac{m_s}{m_{\text{disolución}}} \cdot 100$; $C = \frac{m_s}{V(L)}$ \\
        \textbf{Solución:}
        \begin{enumerate}[label=\alph*)]
        \item \textbf{Porcentaje en masa:}
        Primero calculamos la masa total: $m_{\text{total}} = 1.022 \text{ g/cm}^3 \cdot 200 \text{ cm}^3 = 204.4 \text{ g}$. \\
        $\%m = \frac{12 \text{ g}}{204.4 \text{ g}} \cdot 100 \approx 5.87\%$
        \item \textbf{Concentración en g/L:}
        Dividimos la masa del soluto entre el volumen en litros ($200 \text{ mL} = 0.2 \text{ L}$): \\
        $C = \frac{12 \text{ g}}{0.2 \text{ L}} = 60 \text{ g/L}$
        \end{enumerate}
    \end{solu}
\end{exercise}

\begin{exercise}
    Se puede preparar un suero salino isotónico disolviendo 0,9 g de cloruro de sodio en agua hasta obtener 100 mL de disolución.
    Calcula:
    \begin{enumerate}[label=\letraminuscula]
            \item la concentración en masa en g/L de esta disolución.
            \item la masa de cloruro de sodio necesaria para preparar 250 mL de suero salino isotónico.
    \end{enumerate}
    \begin{solu}
        \textbf{Fórmulas:} $C \text{ (g/L)} = \frac{m_{\text{soluto (g)}}}{V_{\text{disolución (L)}}}$ y $m_s = C \cdot V$ \\
        \textbf{Solución:}
        \begin{enumerate}[label=\alph*)]
                \item \textbf{Concentración:} Cambiamos el volumen a litros ($100 \text{ mL} = 0.1 \text{ L}$). \\
                $C = \frac{0.9 \text{ g}}{0.1 \text{ L}} = 9 \text{ g/L}$.
                \item \textbf{Masa para 250 mL:} Cambiamos el volumen ($250 \text{ mL} = 0.25 \text{ L}$). \\
                $m = 9 \text{ g/L} \cdot 0.25 \text{ L} = 2.25 \text{ g}$ de NaCl.
        \end{enumerate}
    \end{solu}
\end{exercise}


\section{Solubilidad}

\begin{definition}[Solubilidad]
        La \textbf{solubilidad} de un soluto a una cierta temperatura y en un cierto disolvente es la cantidad máxima de ese soluto que puede disolverse en \qty{100}{\g} de disolvente.
\end{definition}

A partir de esta definición podemos ver que la solubilidad depende de la temperatura y del disolvente.
Esto significa que la solubilidad cambia cuando cambia la temperatura y cuando cambia el disolvente.

Por ejemplo, si algo no se disuelve a temperatura ambiente, podemos calentarlo.
Esto demuestra que cuando la temperatura aumenta, la solubilidad también aumenta.

Para expresar correctamente la solubilidad, debemos dar el valor de la solubilidad, indicar que está referida a \qty{100}{\g} del disolvente utilizado, y también dar la temperatura.

\subsection{Curvas de solubilidad}\label{subsec:curvas-de-solubilidad}

Una \textbf{curva de solubilidad} es una gráfica que muestra los datos de solubilidad de una sustancia pura en \qty{100}{\g} de un disolvente a distintas temperaturas.

La Figura \ref{fig:curva_solubilidad}, muestra las curvas de solubilidad de diferentes sustancias puras en agua.

\section*{Resumen}
\begin{itemize}
    \item \textbf{Soluto:} La sustancia que se disuelve y que suele estar presente en menor cantidad en la mezcla.
    \item \textbf{Disolvente:} La sustancia que disuelve al soluto. Suele ser el componente principal (por ejemplo, el agua en disoluciones acuosas).
    \item \textbf{Disolución:} La mezcla homogénea formada por el soluto y el disolvente ($m_{\text{disolución}} = m_{\text{soluto}} + m_{\text{disolvente}}$) o, para el volumen, \( V_{\text{disolución}} = V_{\text{soluto}} + V_{\text{disolvente}} \).
    \item \textbf{Concentración:} Expresa la proporción de soluto en una cierta cantidad de disolución o disolvente.
    \item \textbf{Densidad ($d$):} Relación entre la masa total de la disolución y el volumen que ocupa ($d = m/V$). Se utiliza para convertir volumen en masa y masa en volumen. Es importante no confundir la densidad con la concentración en masa/volumen.
\end{itemize}

\begin{exercise}
    Utiliza las curvas de solubilidad de la última página para responder a estas preguntas.
    \begin{enumerate}[label=\letraminuscula]
            \item ¿Cuál es la solubilidad del KNO3 a 25 °C y a 45 °C?
            \item ¿Qué masa de cristales de \ch{KNO3} se formará si una disolución saturada de 100 g de agua se enfría de 45 °C a 25 °C? ¿Y de 50 °C a 20 °C?
            \item ¿Cuál es la mayor masa de KNO3 que podría disolverse en 1 kg de agua a 50 °C?
            \item ¿Cuánta agua se necesita para disolver 100 g de nitrato de potasio a 45 °C?
            \item ¿A qué temperatura tiene el nitrato de potasio una solubilidad de 20 g por cada 100 g de agua?
    \end{enumerate}
    \begin{solu}
        Aquí va la solución sería otra solución y mejor hacerlo todo.
    \end{solu}
\end{exercise}

\begin{exercise}
    Utiliza las curvas de solubilidad de \ch{KCl} y \ch{NaOH} de la última página para responder a las preguntas.
    \begin{enumerate}[label=\letraminuscula]
            \item ¿La solubilidad de cuál de estas dos sustancias varía menos con la temperatura?
            \item ¿Cuáles son sus solubilidades a \qty{10}{\celsius}?
            \item ¿Cuál es la solubilidad del \ch{KCl} a \qty{30}{\celsius}? ¿Y la del \ch{NaOH} a 80 °C?
            \item ¿Qué masa de \ch{KCl} se formará si una disolución saturada en 100 g de agua se enfría de
            30 °C a 10 °C?
            \item ¿Qué masa de cristales de \ch{NaOH} se formará si una disolución saturada en 100 g de agua se enfría
            de 80 °C a 10 °C?
            \item ¿A qué temperatura tiene el clorato de sodio una solubilidad de 90 g en 100 g de agua?
    \end{enumerate}
    \begin{solu}
        Aquí va la solución
    \end{solu}
\end{exercise}

%\section{Vocabulario}
%
%
%\begin{vocabulary}[Matter]
%Anything that has mass and occupies space. \newline
%\textbf{Translation to Spanish}: Materia \newline
%% Line 14
%\end{vocabulary}
%
%
%
%\begin{vocabulary}[Atom]
%The smallest unit of an element that retains the chemical properties of that element. \newline
%\textbf{Translation to Spanish}: Átomo \newline
%% Line 41
%\end{vocabulary}
%
%
%\begin{vocabulary}[Molecule]
%A group of atoms chemically bonded together. \newline
%\textbf{Translation to Spanish}: Molécula \newline
%% Line 57
%\end{vocabulary}
%
%
%\begin{vocabulary}[Element]
%A pure substance consisting of only one type of atom. \newline
%\textbf{Translation to Spanish}: Elemento \newline
%% Line 76
%\end{vocabulary}
%
%
%\begin{vocabulary}[Compound]
%A substance formed when two or more chemical elements are chemically combined. \newline
%\textbf{Translation to Spanish}: Compuesto \newline
%% Line 94
%\end{vocabulary}
%
%
%\begin{vocabulary}[Mixture]
%A combination of substances that are not chemically bonded and can be separated physically. \newline
%\textbf{Translation to Spanish}: Mezcla \newline
%% Line 128
%\end{vocabulary}
%
%
%\begin{vocabulary}[Density]
%The mass per unit volume of a substance. \newline
%\textbf{Translation to Spanish}: Densidad \newline
%% Line 161
%\end{vocabulary}
%
%
%\begin{vocabulary}[Solution]
%A homogeneous mixture composed of a solute dissolved in a solvent. \newline
%\textbf{Translation to Spanish}: Disolución \newline
%% Line 219
%\end{vocabulary}
%
%
%\begin{vocabulary}[Solute]
%The substance that is dissolved in a solution. \newline
%\textbf{Translation to Spanish}: Soluto \newline
%% Line 231
%\end{vocabulary}
%
%
%\begin{vocabulary}[Solvent]
%The substance in which the solute dissolves to form a solution. \newline
%\textbf{Translation to Spanish}: Disolvente \newline
%% Line 235
%\end{vocabulary}
%
%
%\begin{vocabulary}[Chemical Change]
%A process in which substances are transformed into new substances with different properties. \newline
%\textbf{Translation to Spanish}: Cambio químico \newline
%% Line 307
%\end{vocabulary}
%
%
%\begin{vocabulary}[Physical Change]
%A change that affects the form or state of matter without altering its chemical composition. \newline
%\textbf{Translation to Spanish}: Cambio físico \newline
%% Line 329
%\end{vocabulary}
%
%
%\begin{vocabulary}[State of Matter]
%One of the distinct forms that matter can exist in, such as solid, liquid, gas, or plasma. \newline
%\textbf{Translation to Spanish}: Estado de la materia \newline
%% Line 371
%\end{vocabulary}
%
%
%
%\newpage
%
%\begin{figure}[H]
%    \centering
%    \begin{tikzpicture}
%        \begin{axis}
%            [
%% --- Control del Grid ---
%            grid=both,
%            major grid style={line width=.2pt,draw=gray!70},
%            minor grid style={line width=.1pt,draw=gray!50},
%% ------------------------
%            width=0.9\textwidth,
%            height=0.7\textwidth,
%            axis x line=bottom,
%            axis y line=left,
%            axis line style=thick,
%            xmax=160, xmin=0,
%            ymax=160, ymin=0,
%            xtick={0,20,40,60,80,100, 120, 140, 160},
%            ytick={0,20,40,60,80,100,120,140,160},
%            minor tick num=3, % Esto crea las líneas pequeñas del grid
%            xlabel={Temperatura ( \unit{\celsius})},
%            ylabel={Solubilidad\vspace*{5pt}\\ \( \left(  \dfrac{\unit{g} \text{de soluto}}{\qty{100}{\g} \text{ de agua}} \right) \)},
%            ylabel style={
%                at={(axis description cs:0.1,1.09)},
%                anchor=south east,
%                xshift=-15pt,
%                rotate=-90,
%                align=center,
%                font=\small
%            },
%            xlabel near ticks,
%            ]
%
%% NaOH
%            \addplot[draw, black, smooth, thick] coordinates {(0,15) (25,25) (60,40) (100,50) (160,71)};
%            \node[font=\footnotesize, anchor=west] at (axis cs:140,60) {\ch{NaOH}};
%
%% NaCl
%            \addplot[blue, smooth, thick] coordinates {(0,35.7) (20,35.9) (40,36.2) (60,36.5) (80,36.8) (100,37) (160,38)};
%            \node[blue, font=\footnotesize, anchor=south] at (axis cs:132,39) {\ch{NaCl}};
%
%% KNO3
%            \addplot[red, smooth, thick] coordinates {(0,13.3) (20,31.6) (40,63.9) (60,108) (77,160)};
%            \node[red, font=\footnotesize, anchor=east] at (axis cs:72,145) {\ch{KNO3}};
%
%% KCl
%            \addplot[green!60!black, smooth, thick] coordinates {(0,28.0) (20,34.0) (40,40.5) (60,46.5) (80,52.0) (100,56.5) (160,74)};
%            \node[green!60!black, font=\footnotesize, anchor=north west] at (axis cs:115,71) {\ch{KCl}};
%
%% NH4Cl
%            \addplot[orange, smooth, thick] coordinates {(0,29.0) (20,37.0) (40,45.0) (60,55.0) (80,66.0) (100,76.0) (160,105)};
%            \node[orange, font=\footnotesize, anchor=south west] at (axis cs:111,89) {\ch{NH4Cl}};
%        \end{axis}
%    \end{tikzpicture}
%    \caption{Curva de solubilidad}
%    \label{fig:curva_solubilidad}
%\end{figure}
%
%
